\documentclass{article}
\usepackage{amsmath, amssymb, amsthm}
\usepackage{hyperref}
\usepackage{geometry}
\geometry{margin=1in}

\title{Erd\H{o}s Problem \#295}
\date{Last updated: 2025-08-31}
\author{Source: erdosproblems.com}

\begin{document}
\maketitle

\noindent\textbf{Status:} OPEN \quad \textbf{No prize}

\noindent\textbf{Tags:} number theory, unit fractions

\medskip

\noindent\textbf{Problem Statement:}

Let $N\geq 1$ and let $k(N)$ denote the smallest $k$ such that there exist $N\leq n_1&#60;\cdots &#60;n_k$ with\[1=\frac{1}{n_1}+\cdots+\frac{1}{n_k}.\]Is it true that\[\lim_{N\to \infty} k(N)-(e-1)N=\infty?\]




Erd\H{o}s and Straus \cite{ErSt71b} have proved the existence of some constant $c&#62;0$ such that\[-c &#60; k(N)-(e-1)N \ll \frac{N}{\log N}.\]




References

[ErSt71b] Erd\H{o}s, P. and Straus, E. G., Solution to Problem . Amer. Math. Monthly (1971), 302-303.

\noindent\textbf{Related OEIS sequences:} A192881


\bigskip
\noindent\small{Source: \url{https://www.erdosproblems.com/295}}
\end{document}
